\documentclass[a4paper,12pt]{article}
\usepackage[utf8]{inputenc}
\usepackage[T1]{fontenc}
\usepackage[ngerman]{babel}
\usepackage{amsmath}
\usepackage{amssymb}
\usepackage{enumitem}
\usepackage{geometry}
\geometry{a4paper, margin=2.5cm}
\usepackage{parskip}

\title{Einsendeaufgaben -- Aufgabe 5.2\\[0.5em]
\large Modul 61112: Lineare Algebra}
\author{}
\date{}

\begin{document}
\maketitle

\section*{Aufgabe 5.2}

\begin{align*}
\chi_A &= \det(XI_4 - A) \\
       &= \det\begin{pmatrix}
            X-\overline{3} & \overline{0}   & \overline{0}   & \overline{0} \\
            \overline{1}   & X-\overline{3} & \overline{2}   & \overline{0} \\
            \overline{0}   & \overline{0}   & X-\overline{1} & \overline{0} \\
            \overline{0}   & \overline{0}   & \overline{0}   & X-\overline{1}
          \end{pmatrix} \\
       &= (X-\overline{3}) \det\begin{pmatrix}
            X-\overline{3} & \overline{2}   & \overline{0} \\
            \overline{0}   & X-\overline{1} & \overline{0} \\
            \overline{0}   & \overline{0}   & X-\overline{1}
          \end{pmatrix} \\
       &= (X-\overline{3})^{2} \det\begin{pmatrix}
            X-\overline{1} & \overline{0} \\
            \overline{0}   & X-\overline{1}
          \end{pmatrix} \\
       &= (X-\overline{3})^{2}\, (X-\overline{1})^{2}
\end{align*}

Das Minimalpolynom $\mu_A$ teilt nach Satz von Cayley-Hamilton das charakteristische Polynom $\chi_A$. Da $\chi_A$ in Linearfaktoren zerfällt, also $\chi_A = (X-\overline{3})^{2}\, (X-\overline{1})^{2}$, muss nach Lemma 1.4.37 $\mu_A$ auch in diese Linearfaktoren zerfallen.

$(X-\overline{3})$ und $(X-\overline{1})$ liegen nicht in $I_A$. Es existiert also kein Teiler vom Grad 1 in $I_A$.

Wir testen die Teiler vom Grad 2:

\[
(A - \overline{3}I_4)^2 = \begin{pmatrix}
  \overline{0} & \overline{0} & \overline{0}  & \overline{0} \\
  \overline{1} & \overline{0} & \overline{2}  & \overline{0} \\
  \overline{0} & \overline{0} & \overline{-2} & \overline{0} \\
  \overline{0} & \overline{0} & \overline{0}  & \overline{-2}
\end{pmatrix}^{\!2}
= \begin{pmatrix}
  \overline{0} & \overline{0} & \overline{0}  & \overline{0} \\
  \overline{0} & \overline{0} & \overline{-4} & \overline{0} \\
  \overline{0} & \overline{0} & \overline{4}  & \overline{0} \\
  \overline{0} & \overline{0} & \overline{0}  & \overline{4}
\end{pmatrix}
\]

\[
(A - \overline{1}I_4)^2 = \begin{pmatrix}
  \overline{2} & \overline{0} & \overline{0} & \overline{0} \\
  \overline{1} & \overline{2} & \overline{2} & \overline{0} \\
  \overline{0} & \overline{0} & \overline{0} & \overline{0} \\
  \overline{0} & \overline{0} & \overline{0} & \overline{0}
\end{pmatrix}^{\!2}
= \begin{pmatrix}
  \overline{4} & \overline{0} & \overline{0} & \overline{0} \\
  \overline{4} & \overline{4} & \overline{4} & \overline{0} \\
  \overline{0} & \overline{0} & \overline{0} & \overline{0} \\
  \overline{0} & \overline{0} & \overline{0} & \overline{0}
\end{pmatrix}
\]

\[
(A - \overline{3}I_4)(A - \overline{1}I_4) =
\begin{pmatrix}
  \overline{0} & \overline{0} & \overline{0}  & \overline{0} \\
  \overline{1} & \overline{0} & \overline{2}  & \overline{0} \\
  \overline{0} & \overline{0} & \overline{-2} & \overline{0} \\
  \overline{0} & \overline{0} & \overline{0}  & \overline{-2}
\end{pmatrix}
\begin{pmatrix}
  \overline{2} & \overline{0} & \overline{0} & \overline{0} \\
  \overline{1} & \overline{2} & \overline{2} & \overline{0} \\
  \overline{0} & \overline{0} & \overline{0} & \overline{0} \\
  \overline{0} & \overline{0} & \overline{0} & \overline{0}
\end{pmatrix}
= \begin{pmatrix}
  \overline{2} & \cdot & \cdot & \cdot \\
  \cdot & \cdot & \cdot & \cdot \\
  \cdot & \cdot & \cdot & \cdot \\
  \cdot & \cdot & \cdot & \cdot
\end{pmatrix}
\]

Dennoch liegt auch kein Teiler vom Grad 2 in $I_A$.

Aber $g \in \mathbb{F}_p[X]$ mit $g := (X-\overline{3})^2 (X-\overline{1})$ liegt in $I_A$, da
\begin{align*}
g(A) &= (A - \overline{3}I_n)^2 (A - \overline{1}I_n) \\
     &= \begin{pmatrix}
       \overline{0} & \overline{0} & \overline{0}  & \overline{0} \\
       \overline{0} & \overline{0} & \overline{-4} & \overline{0} \\
       \overline{0} & \overline{0} & \overline{4}  & \overline{0} \\
       \overline{0} & \overline{0} & \overline{0}  & \overline{4}
     \end{pmatrix}
     \begin{pmatrix}
       \overline{2} & \overline{0} & \overline{0} & \overline{0} \\
       \overline{1} & \overline{2} & \overline{2} & \overline{0} \\
       \overline{0} & \overline{0} & \overline{0} & \overline{0} \\
       \overline{0} & \overline{0} & \overline{0} & \overline{0}
     \end{pmatrix} \\
     &= 0
\end{align*}

Da $\mu_A$ aufgrund der vorherigen Überlegungen vom Grad $\geq 3$ sein muss, ist $\mu_A = g = (X-\overline{3})^2 (X-\overline{1})$.

\end{document}
